% August 2012
% Autor: Michael Meese & Thomas Fromme
% Universität Freiburg 
%www.fh-bielefeld.de

\documentclass{beamer} 								% Standardklasse für Beamer
	%\documentclass[handout]{beamer}  			% druckoptimiert (handouts)
   	%\documentclass[draft]{beamer}					% Entwickler-Modus in dem das File schneller kompiliert wird
	%\documentclass{article}							% Artikel-Version; Funktioniert nur, wenn man zusätzlich das Paket beamerarticle mittels \usepackage{beamerarticle} einbindet.

\usetheme{progressbar}
\usepackage[ngerman]{babel}
\usepackage[utf8x]{inputenc}
	%\usepackage[latin1]{inputenc}						% Für Windows
	% \usepackage[ansinew]{inputenc} 				% Für Windows

\usepackage{amsmath,amsfonts,amssymb}
	\progressbaroptions{imagename=fhlogo}
	\usecolortheme{default}							% albatross | beaver | beetle | crane | default | dolphin | dove | fly | lily | orchid | rose |seagull | seahorse | sidebartab | structure | whale | wolverine
	\usefonttheme{professionalfonts}	
															% default | professionalfonts | serif | structurebold | structureitalicserif | structuresmallcapsserif
	\useinnertheme{inmargin}
															% circles | default | inmargin | rectangles | rounded
	\useoutertheme{smoothbars}
															% default | infolines | miniframes | shadow | sidebar | smoothbars | smoothtree | split | tree
	%\setbeamercovered{transparent}				% halbtransparente Overlays auf seinen Folien
	%\beamertemplatenavigationsymbolsempty	% abschalten der Navigationselemente
	%\setbeamertemplate{footline}[frame number]	% Seitenzahlen in Fußzeile einfügen


\begin{document}
	\title[Abbrev. Title\hspace{2em}\insertframenumber/
		\inserttotalframenumber]{Stereoskopie}
	\author{Michael Meese & Thomas Fromme} 

	\date{\today}
	%\logo{\includegraphics[scale=0.14]fhlogo}}

	
	%\title[Kurzform]{Vortrag zur Berechenbarkeit}		% Titel des Vortrages
	%\subtitle[Kurzform]{Untertitel}											% Untertitel
	 \author[T.Fromme]{Thomas Fromme}							% Autor festlegen
	 \institute[FB2 -- FH-Bielefeld]{University of Applied Sciences}
																													% Angabe des Institutes

	%\date[26.05.06]{26. Mai 2006}											%Datum der Präsentation, alternativ kann mittels \date{\today} auch das aktuelle Datum eingetragen werden.
\logo{\pgfimage[width=2cm,height=2cm]{img/fhlogo}}
  																%Die Datei hulogo.pdf (bzw. hulogo.png, hulogo.jpg, hulogo.mps bei Verwendung von 
																%pdftex als Backend) als Logo auf allen Folien, hier mithilfe des Paketes pgf.
\titlegraphic{\includegraphics[width=2cm,height=2cm]{img/fhlogo}}
																%Die Datei hulogo.pdf (bzw. analog wie bei \logo auch entsprechendes Format) als Logo
																%nur auf der Titelseite unter Verwendung des Paketes graphicx.
\subject{Stereoskopie}									%Setzt Thema in den PDF-Dokument-Eigenschaften
\keywords{Stereoskopie, OpenGL, OpenCL}				%Setzt Schlüsselwörter in den PDF-Dokument-Eigenschaften

\begin{frame}
	\titlepage
\end{frame}

\begin{frame}
	\frametitle{Inhaltsverzeichnis}\tableofcontents
\end{frame}
%##########################################################################################%
%##########################################################################################%
%##########################################################################################%
%\part[Kurzform]{Name}					%Bei Lesungen mit mehreren Vorträgen bietet sich eine Unterteilung in Parts an, da dann jeder Vortrag mit eigenem Inhaltsverzeichnis und Navigationsstrukturen aufwartet.
%\section[Kurzform]{Name}					%Erzeugt eine neuen Abschnitt mit entsprechendem Namen
%\subsection[Kurzform]{Name}				%Erzeugt einen neuen Unterabschnitt mit entsprechendem Namen 
%##########################################################################################%
%##########################################################################################%
%##########################################################################################%















\section[Quellen]{Referezen}
\begin{frame}\frametitle{Quellen \& Literatur}
%\begin{thebibliography}{1}
%\bibitem[Stereoskopie]{bpt} \emph{Stereoskopie} \\
%\url{http://stereoskopie.com} [Abruf: 16.03.2010]\\
%\end{thebibliography}

\begin{thebibliography}{1}
\bibitem[ ]{bpt} \emph{ } \\
\url{http:// } [Abruf: 20.07.2012]\\
\end{thebibliography}

\begin{thebibliography}{2}
\bibitem[ ]{bpt} \emph{ } \\
\url{http:// } [Abruf: 20.07.2012]\\
\end{thebibliography}

\end{frame}
\end{document}
